\section{Discussion}
\label{sec:discussion}

\subsection{Meaning and limits of the PA target interval}

The results define a conditional PA prototype target interval without
establishing a general PA thermal advantage. In the small-diameter, thin-wall
regime, the outer convective resistance dominates, so the low PA wall conductivity affects
$k_{\mathrm{o}}$ mainly through the wall-resistance share
$\phi_{\mathrm{w}}$. This is consistent with the near-matched
PA612--stainless-steel comparison~\cite{Astrouski2026}, where a staggered
$d_{\mathrm{o}}=1.23\,\mathrm{mm}$, $t=0.12\,\mathrm{mm}$ PA612 bank at
$2$--$10\,\mathrm{m\,s^{-1}}$ lost ${\approx}\,7\,\%$ of mean heat rate; its
staggered layout, thicker relative wall, and steel comparator test the
resistance-share mechanism for a different geometry, whereas the present map
uses inline benchmark pitch and the combined constraints.
Positive values in Fig.~\ref{fig:tech_adj} mean that a feasible PA
point exceeds the nearest feasible aluminum alternative under the constraints
used in Fig.~\ref{fig:design_boundary}.
The PA interval depends on more than pure area scaling; it emerges from the
combined effect of outer-area conductance, resistance shares, the fixed
packing assumptions, and the applied tolerance, potting, burst, and cost
limits.

This interval is bounded by manufacturing and durability constraints. The PA
region depends on wall thicknesses above the nominal process minimum under the
normal tolerance case, controlled potting, and the derated burst-strength
input.
Although the lower PA density favors tube mass, assembly-level mass and
performance can still be dominated by potting, supports, tube sheets, bypass
flow, and pressure integrity.

\Needspace{15\baselineskip}
\subsection{Manufacturing and durability constraints}

Wall tolerance is the first practical limit. At
$d_{\mathrm{o}}=0.5\,\mathrm{mm}$ and $t=0.025\,\mathrm{mm}$, an eccentricity of
$\pm 5\,\mu\mathrm{m}$ is already a $\pm 20\,\%$ wall-thickness variation.
For the normal route used in the integrated map,
$\Delta t=20\,\mu\mathrm{m}$ leaves a
$5\,\mu\mathrm{m}$ local wall at the PA process minimum; the burst
criterion forces walls above $t_{\mathrm{min},\mathrm{PA}}$. A Lam\'e pressure
check at nominal wall thickness would be non-conservative. Thin-walled
sub-millimeter PA hollow fibers have been demonstrated in heat-exchanger
applications~\cite{Kroulikova2021,Bulejko2022}, but tolerances, concentricity,
continuous length, tube-count yield, and cost at a given diameter--wall
combination remain process- and supplier-specific~\cite{Mount2017}; public
supplier capability data do not document the full combination assumed here.

Potting is the second practical limit. The capillary-rise limit in
Fig.~\ref{fig:design_boundary} is a potting-process constraint: resin rise
shortens active length but does not change the local heat-transfer
correlations. Centrifugal potting addresses the body-force term in the
capillary balance, so membrane-module practice gives useful process guidance.
For this heat-exchanger geometry, however, the $G=10$ assumption still has to be
checked with the final tube surface, resin system, gel time, fixture, and
potting depth.

Pressure integrity must be assessed at the supported potted bundle level;
isolated tube-wall burst covers a single failure mode. In early prototype
builds, failures occurred in the frame, potting, and tube--potting interface,
with bundle deformation contributing to the failure path. The $6\,\mathrm{bar}$ Lam\'e
criterion excludes tube-wall geometries with inadequate idealized pressure
margin. Supports, creep, aging, fatigue, pressure cycling, moisture uptake and
hydrolysis in the glycol--water coolant, and the adhesion and long-term
durability of the tube--potting interface are bundle qualification items.
Published pulsating-pressure fatigue
tests~\cite{Brozova2016} and temperature-dependent burst
measurements~\cite{Bulejko2022} of polymeric hollow-fiber heat-transfer
surfaces provide first experimental anchors for these mechanisms; the present
bundle still needs its own support, potting, and coolant-exposure tests.
Higher-temperature
polymers such as polyether ether
ketone (PEEK) may improve these margins but require their own material/process
sweep and tests~\cite{Peters2017}.

\subsection{Model scope}

The heat-transfer model is a VDI-correlation sweep at one operating point with
constant fluid properties and benchmark inline pitch
ratios~\cite{GnielinskiG7,GnielinskiG1}; the resulting
$k_{\mathrm{o}}A_{\mathrm{o}}$ gives a package-scaled conductance potential.
Full heat-rate and vehicle-system predictions require heat-rate and
outer- and tube-side pressure-drop evaluation including
headers and bypass, fan and pump power, and tests on the realized bundle. The
capillary map informs potting trials, and the tube-supply-cost index flags
diameter ranges for supplier quotation rather than setting a production cost.
